\begin{frame}{工作负载与指标的先后关系}
  \Core{指标不能脱离场景单独解释；只有先交代 workload，数字才具有明确含义。}
  \begin{itemize}
    \item 离线批处理更接近资源利用率问题，在线对话更接近响应性问题。
    \item RAG 与 Agent 会把“生成”嵌入更长链路，因此单点吞吐不再足以代表体验。
    \item 同一套指标在不同 workload 下会具有不同含义。
  \end{itemize}
\end{frame}

\begin{frame}{例子：不同在线 workload 的差异}
  \Core{即便都属于在线请求，不同 workload 的长度分布、链路结构与服务目标也可能完全不同。}
  \small
  \begin{tabular}{p{0.20\textwidth}p{0.29\textwidth}p{0.31\textwidth}}
    \toprule
    场景 & 分布特征 & 更敏感的系统问题 \\
    \midrule
    对话服务 & 到达连续，prompt 长度波动大 & TTFT、尾时延、短请求体验 \\
    RAG 服务 & 检索与生成交织，链路更长 & 全链路时延、阶段协同 \\
    Agent 服务 & 多轮工具调用与多次生成 & 任务级完成时间、跨轮状态维护与一致性 \\
    \bottomrule
  \end{tabular}
\end{frame}

\begin{frame}{例子：均值为何经常产生误导}
  \Core{系统实验中最常见的误区，是只看均值而忽视分布、尾部与服务目标。}
  \begin{tabular}{p{0.26\textwidth}p{0.28\textwidth}p{0.28\textwidth}}
    \toprule
    表面现象 & 可能遗漏的信息 & 更合适的补充指标 \\
    \midrule
    平均 TTFT 正常 & 长请求与突发到达已推高 P99 & P95/P99、队列等待 \\
    平均吞吐更高 & 短请求体验可能显著恶化 & TTFT、短请求完成时间 \\
    完成量增加 & 满足 SLO 的请求比例可能下降 & Goodput、SLO 达成率 \\
    \bottomrule
  \end{tabular}
\end{frame}

\begin{frame}{代码例子：benchmark 如何固定分布并生成指标}
  \Core{工作负载与指标分析需要学生看到：实验结果并非“随便跑一遍”即可解释。}
  \begin{columns}[T,totalwidth=\textwidth]
    \column{0.5\textwidth}
    \begin{block}{输入分布控制}
      {\ttfamily\small
      \detokenize{seed(0)}\par
      \detokenize{prompt_token_ids = [[randint(...) ...]]}\par
      \detokenize{sampling_params = [SamplingParams(...)]}\par
      }
    \end{block}
    \column{0.46\textwidth}
    \begin{block}{结果计算逻辑}
      {\ttfamily\small
      \detokenize{t = time.time()}\par
      \detokenize{llm.generate(..., use_tqdm=False)}\par
      \detokenize{throughput = total_tokens / (time.time() - t)}\par
      }
    \end{block}
  \end{columns}
  \vspace{0.4em}
  \begin{itemize}
    \item 输入长度分布与输出长度分布需要分别说明。
    \item 仅有总体 throughput 仍不足以支撑 TTFT、P99 与队列等待层面的完整结论。
  \end{itemize}
\end{frame}

\begin{frame}{例子：throughput、TTFT 与 P99 的张力}
  \Core{工作负载与评价指标这一讲的重点，是接受真实系统中的指标张力并据此陈述取舍。}
  \begin{itemize}
    \item 更大的 batch 往往提高总体吞吐，但也可能抬高等待时间与 TTFT。
    \item 更保守的调度可能稳定尾时延，却压低设备利用率与平均吞吐。
    \item 因而任何优化都应首先说明优先保留的指标与可接受的代价。
  \end{itemize}
\end{frame}

\begin{frame}{课堂练习：为同一系统写出不同的指标解释}
  \Core{同一组数字可以对应不同解释，关键在于问题定义、工作负载设定与目标指标选择。}
  \begin{enumerate}
    \item 若一个系统吞吐提高 15\%，但 TTFT 恶化 30\%，如何判断它是否更适合在线服务。
    \item 给出一个适用于 Agent 工作流的 SLO 组合，而不只写单个 TTFT 阈值。
    \item 说明为什么指标解释必须回到 workload 与场景。
  \end{enumerate}
\end{frame}
